% page 574 of the facsimile (image p-573 of the scan)
% block 152.4 x 254.9 mm ; left margin 8.0 mm
\pgc{-12.700mm}{0.000mm}{
\l{}
\l{}
\l{}
\l{}
\l{\cel{4}566}
\l{}
\l{\cel{1}\sou{tabureto}. I. Sidilo kun quar pedi, sen brakii nek dors-}
\l{\cel{4}apogilo. - II. Mikra suportilo, di qua la formo esas}
\l{\cel{4}analoga, sur qua onu pozas sua pedi kande onu sidas}
\l{\cel{4}sur stulo. - DFRS.}
\l{}
\l{\cel{1}tacar. (netrans., pri). I.Silencar, neparolar, pri (ulo,}
\l{\cel{4}ulu). - II. Permanar silencoza. - EFIS.}
\l{}
\l{\cel{1}\sou{tafto}. Stofo ek silko, uniforma e briloza. - DEFIRS.}
\l{}
\l{\cel{1}tago. Peceto de metalo, ye la extremajo di laco, di kor-}
\l{\cel{4}dono. - E.}
\l{}
\l{\cel{1}tako. (tekn.) Mikra peco de ligno qua subtenas la extrema-}
\l{\cel{4}jo di traverseto ; qua mantenas angulo di moblo ; qua}
\l{\cel{4}garnisas la kuseneto di pulio por ke olca ne tendencez}
\l{\cel{4}a sterno ; quan onu sinkas aden tero por uzesar kom}
\l{\cel{4}repero-punto lor la rektigo di voyo ; qua startas la}
\l{\cel{4}naveto \textquotedbl{}fluganta\textquotedbl{} di texo-mashino, e c. - FI.}
\l{}
\l{\cel{1}\sou{takigenezo}. (\sou{biol.)}\cel{2}Formo di la kresko en qua ula fazi}
\l{\cel{4}embrionifala esas kondensita o supresita.\cel{2}DEFIS.}
\l{}
\l{+\sou{takigrafar}. (\sou{trans.}) I. (en\cel{1}\sou{la epoki antiqua,}\cel{1}en Roma).}
\l{\cel{4}Uzar, por la enrejistro di lo dicata, la sistemo de}
\l{\cel{4}stenografo-signi quan inventis Tiro, la sekretario di}
\l{\cel{4}Markus Tulius Kikero (Cicero). - II. (nia-epoke)}
\l{\cel{4}Mi-stenografar, t.e. uzar abrevio-sistemo qua posiblig-}
\l{\cel{4}as notar ye rapideso-grado pasable plu granda kam per}
\l{\cel{4}la skribo-sistemo normala, ma ne tam rapide kam per ste-}
\l{\cel{4}nografo. - DEFIRS.}
\l{}
\l{\cel{1}\sou{takimetro}. Aparato uzata pri mekaniko, por mezurar la rapi-}
\l{\cel{4}deso-gradi angulala, o grandeso irg-altra qua esas fun-}
\l{\cel{4}ciono lineara di ta rapideso. - DEFIRS.}
\l{}
\l{\cel{1}\sou{takimetrio}. Metodo qua posibligas demonstrar rapide la}
\l{\cel{4}teoremi geometriala od algebrala precipua, per konkret-}
\l{\cel{4}igar la operaci quin eventigas la demonstro. - DEFIRS.}
\l{}
\l{\cel{1}takto. Sento delikata pri to quo konvenas agar o ne agar}
\l{\cel{4}por ne shokar o shameto-redigar la personi altra. -}
\l{\cel{3}DEFIRS.}
\l{}
\l{\cel{1}taktiko. I. (armeo).La arto situar e movigar, sur tereno,}
\l{\cel{4}la trupi di armo-speci diversa. - II. (\sou{metaf.}) La ma-}
\l{\cel{4}niero ye qua onu agas por atingar ta o ca skopo, por ob-}
\l{\cel{4}tenar ta o ca rezultajo. - DEFIRS.}
\l{}
\l{\cel{1}\sou{taktismo.(biol.}) Influo quan agas ula substanci kemiala od}
\l{\cel{4}ula formi di energio (koloro, kaloro, foto, elektro) a}
\l{\cel{4}la protoplasmo e precipue a la direciono, la translaco}
\l{\cel{4}di la celuli libera. - DEFIS.}
}
